% latex-oef2.tex
%
\documentclass[11pt,a4paper]{amsart}
\usepackage{amsmath,amssymb,amsthm}
\usepackage[utf8]{inputenc}
\usepackage[dutch]{babel}
\usepackage{enumitem}
\usepackage{graphicx}
\begin{document}
Combinatoriek
Variaties
We trekken k keer uit een verzameling van n elementen, zonder 
terugleggen en onthouden de volgorde. Voor de eerste trekking zijn er 
n 
mogelijkheden. Voor elke volgende trekking is er steeds één mogelijkheid minder 
dan voor de vorige. Een getrokken rijtje heet een variatie. Daarvan 
zijn er
Voorbeeld: vier vrienden hebben iets te vieren en ze hebben een complete zaal 
gehuurd. Er zijn 200 plaatsen en ze gaan zo maar ergens zitten. Hoeveel 
mogelijkheden zijn er?
Als we alle n elementen trekken, dus als k = n, krijgen we een 
mogelijke manier om de n elementen van de verzameling op volgorde te zetten. 
Zo'n volgorde heet een permutatie. Daarvan zijn er
Voorbeeld: we sorteren een pak kaarten (trekken 52 kaarten uit het pak); er
zijn dan 52! mogelijkheden.

Herhalingsvariaties
We trekken k keer uit een verzameling van n elementen, met 
terugleggen en onthouden de volgorde. Voor elk van de k trekkingen 
zijn er n mogelijkheden. Een getrokken rijtje heet een 
herhalingsvariatie. Daarvan zijn er
Voorbeeld: Hoeveel pincodes zijn er? Men kiest 4 maal uit 10 cijfers dus .

Combinaties
We trekken k keer uit een verzameling van n elementen, zonder 
terugleggen en letten niet op de volgorde. Een getrokken k-tal heet 
een 
combinatie. We kunnen eerst wel op de volgorde letten en die daarna 
vergeten. Er zijn dan k! variaties die hetzelfde k-tal opleveren. Het 
aantal 
verschillende combinaties is dus:
Voorbeeld: Hoeveel verschillende ploegen van 3 deelnemers kunnen we vormen uit 
10 deelnemers. Dat zijn er

Herhalingscombinaties
We trekken k keer uit een verzameling van n elementen, met 
terugleggen 
en letten niet op de volgorde. Een getrokken k-tal heet een 
herhalingscombinatie. We kunnen zo'n k-tal beschrijven door van elk van de 
n elementen aan te geven hoe vaak het gekozen is. Dit kan aanschouwelijk 
gebeuren door voor elk element net zoveel nullen te schrijven als het gekozen 
is 
en tussen de nullen van de verschillende elementen een 1 als scheiding te 
schrijven. Het aantal verschillende herhalingscombinaties is dan het aantal 
mogelijke rijtjes bestaande uit k keer een 0 en n-1 keer een 1. We 
moeten van 
de k+n-1 plaatsen in de rij er k aanwijzen waar een 0 komt te staan. 
Het aantal 
is dus net zoveel als het aantal combinaties van k uit k+n-1:
Voorbeeld: Op hoeveel manieren kun je 5 eieren kleuren als je over 3 kleuren 
beschikt. Dus k=5 en n=3. Dat kan dus op
manieren.
Toepassing
Op hoeveel manieren kunnen n gelijke knikkers verdeeld worden over k 
verschillende bakjes, zo, dat geen bakje leeg is?

Doe in elk bakje alvast een knikker. Er blijven er n-k over. Deze worden 
nog over de bakjes verdeeld via een herhalingscombinatie, dus van n-k uit 
k; voor elk van de resterende knikkers wordt bepaald in welk bakje hij 
komt. Er zijn dus:
mogelijkheden.
\end{document}
